\chapter{Conclusions and Recommendations}
\label{ch:conclusions}

This chapter synthesises the findings of the thesis, provides
direct answers to the research questions stated in
Chapter~\ref{ch:introduction}, summarises the contributions of
the work, presents specific recommendations for the Nepal Rastra
Bank, and identifies directions for future research.

% ================================================================
\section{Summary of Work}
\label{sec:summary}
% ================================================================

This thesis developed a formal, normalised metadata
de-anonymization risk model for token-based Central Bank Digital
Currency and applied it to Nepal's domestic CBDC deployment
context under Nepal Rastra Bank jurisdiction.

The theoretical framework derived three interconnected components
from first principles. The single-attribute uniqueness score
$s(a_i)$ measures the fraction of users uniquely identifiable
from a single metadata attribute. The pairwise interaction
coefficient $\lambda_{ij}$ captures the amplification in
re-identification risk when two attributes are observed in
combination, beyond what either attribute reveals individually.
The normalised risk score $R_{\mathrm{norm}} = \Rraw / \Rmax$
maps total accumulated risk into the unit interval using a proven
theoretical upper bound $\Rmax(m) = m + \binom{m}{2}$, enabling
direct comparison across schemes, attribute sets, and population
sizes.

The model was applied to six cryptographic schemes --- Plain Hash
Token, Blind Signature, Hybrid Scheme, Ring Signature, Stealth
Address, and ZKP Token --- across seven metadata attributes
including two Nepal-specific fields: payment channel and
transaction zone. Eight analyses were conducted covering baseline
scheme comparison, population scalability, mitigation
effectiveness, attribute sensitivity, seed sensitivity,
micro-scale pilot behaviour, lambda sensitivity, and worst-case
calibration against the theoretical maximum.

% ================================================================
\section{Answers to Research Questions}
\label{sec:answers}
% ================================================================

\textbf{Primary question:} How can the metadata de-anonymization
risk of different cryptographic schemes in a token-based CBDC be
formally measured, normalised, and compared on a common
quantitative scale?

The $\Rnorm$ metric answers this question directly. By
normalising $\Rraw$ against its proven theoretical maximum
$\Rmax(m)$, the metric produces a scale-independent score in
$[0, 1]$ comparable across any number of attributes and any
population size. The metric was validated across 30 random seeds
and calibrated against a worst-case scenario achieving
$\Rnorm = 0.9994$, confirming the normalisation is correctly
calibrated at 99.9\% of the theoretical maximum.

\textbf{Secondary question 1 --- Scheme comparison:}
At baseline ($N = 1{,}000$, $m = 7$), the six schemes produce
the following $\Rnorm$ values: Plain Hash Token 0.317, Blind
Signature 0.174, Stealth Address 0.107, Ring Signature 0.069,
Hybrid Scheme 0.067, and ZKP Token 0.035. The six schemes form
a clean, non-overlapping risk ladder. ZKP Token is 9.1 times
safer than Plain Hash Token. Hybrid Scheme and Ring Signature
are statistically equivalent at $N = 1{,}000$.

\textbf{Secondary question 2 --- Population scalability:}
Scale does not resolve the privacy problem for weak schemes.
Plain Hash Token at $N = 100{,}000$ carries $\Rnorm = 0.191$
--- 5.5 times the risk of ZKP Token at the same scale. ZKP Token
is the only scheme whose $\Rnorm$ is scale-independent, remaining
flat at 0.035 from $N = 100$ to $N = 100{,}000$. This formally
disproves the assumption that population growth can substitute
for cryptographic privacy design.

\textbf{Secondary question 3 --- Mitigation effectiveness:}
A full mitigation stack reduces risk substantially across all
schemes. Ring Signature with full mitigation reaches
$\Rnorm = 0.002$, within 0.001 of ZKP Token with full mitigation.
Wallet unlinkability is the single most effective individual
mitigation across all schemes. Nepal Controls alone produce zero
risk reduction at $N = 1{,}000$ because Nepal's demographic
clustering already suppresses uniqueness on payment channel and
transaction zone.

\textbf{Secondary question 4 --- Nepal-specific context:}
Adding Nepal-specific attributes counterintuitively
\textit{decreases} $\Rnorm$ for all schemes because $\Rmax$
grows by 6 then 7 as each attribute is added, while $\Rraw$
increases by near-zero due to the skewed distributions of these
attributes. This formally proves that Nepal's demographic
structure provides passive privacy protection on payment channel
and transaction zone fields that does not require explicit
technical mitigation at domestic scale.

% ================================================================
\section{Contributions}
\label{sec:contributions}
% ================================================================

This thesis makes five contributions to the existing body of
knowledge.

\textbf{Contribution 1 --- Formal normalised privacy risk metric.}
The $\Rnorm$ metric provides the first formally derived,
normalised privacy risk score for token-based CBDC metadata.
Unlike k-anonymity and differential privacy, which were not
designed for scheme comparison, $\Rnorm$ produces a continuous
score in $[0, 1]$ directly comparable across any cryptographic
scheme, attribute set, and population size. The derivation of
$\Rmax(m) = m + \binom{m}{2}$ as the proven theoretical upper
bound is an original mathematical contribution.

\textbf{Contribution 2 --- Quantitative six-scheme comparison.}
This thesis provides the first quantitative comparison of six
CBDC cryptographic schemes on a common privacy risk metric. Prior
to this work, inter-scheme comparisons were qualitative. The risk
ladder spanning $\Rnorm = 0.317$ to $\Rnorm = 0.035$ gives
policymakers specific numbers for the privacy cost of each design
choice.

\textbf{Contribution 3 --- Formal disproof of scale substitution.}
The population scalability analysis formally disproves the
assumption that deploying a CBDC across a larger population
provides meaningful privacy improvement for weak schemes. Plain
Hash Token's residual risk floor of 0.191 at $N = 100{,}000$
demonstrates that cryptographic design, not population size, is
the determinant of long-run privacy risk.

\textbf{Contribution 4 --- Nepal-specific deployment analysis.}
By incorporating payment channel and transaction zone as
Nepal-specific attributes grounded in Nepal's payment
infrastructure, this thesis provides the first formal privacy risk
analysis specifically designed for Nepal's domestic CBDC context.

\textbf{Contribution 5 --- Micro-scale pilot risk quantification.}
The micro-scale analysis demonstrates that at VDC pilot scale
($N = 50$), all schemes except ZKP Token exceed
$\Rnorm = 0.48$. This finding directly informs how NRB should
approach any community-level CBDC pilot.

% ================================================================
\section{Recommendations for Nepal Rastra Bank}
\label{sec:nrb_recommendations}
% ================================================================

Based on the findings of this thesis, the following specific
recommendations are made to the Nepal Rastra Bank.

\textbf{Recommendation 1 --- Do not deploy Plain Hash Token at
any scale.} Plain Hash Token produces $\Rnorm = 0.317$ at pilot
scale and 0.191 at national scale. Both values represent
unacceptable citizen re-identification risk under Nepal's
constitutional right to privacy. The scheme should be excluded
from all deployment consideration.

\textbf{Recommendation 2 --- Adopt Ring Signature as the minimum
acceptable scheme for national deployment.} Ring Signature
achieves $\Rnorm = 0.069$ at baseline and 0.002 with a full
mitigation stack --- near-ZKP performance at significantly lower
computational cost.

\textbf{Recommendation 3 --- Consider Hybrid Scheme for
resource-constrained deployment.} Hybrid Scheme
($\Rnorm = 0.067$) is statistically equivalent to Ring Signature
($\Rnorm = 0.069$) at $N = 1{,}000$ and may be implementable at
lower computational cost on Nepal's low-end rural devices.

\textbf{Recommendation 4 --- Prioritise wallet unlinkability in
any mitigation stack.} Wallet unlinkability produces the largest
single reduction in $\Rnorm$ across all schemes. Any deployed
CBDC wallet should generate fresh addresses per transaction or
implement nullifier-based unlinkability.

\textbf{Recommendation 5 --- Do not rely on Nepal Controls as a
substitute for cryptographic privacy.} Nepal Controls produce
zero risk reduction at $N = 1{,}000$. Mitigation resources should
be focused on amount, timestamp, and wallet attributes where
marginal reduction is largest.

\textbf{Recommendation 6 --- Use ZKP Token only if computational
constraints can be resolved.} ZKP Token is the only scheme with
scale-independent risk. If NRB can resolve its computational
requirements through tiered CBDC architecture or offline ZKP
verification, it should be the target deployment scheme.

\textbf{Recommendation 7 --- Conduct any village pilot under
Ring Signature or better.} At $N = 50$, all schemes except ZKP
Token exceed $\Rnorm = 0.48$. Any NRB community-level pilot must
deploy Ring Signature or above to protect pilot participants from
high re-identification risk during the experimental phase.

% ================================================================
\section{Limitations}
\label{sec:limitations_ch5}
% ================================================================

The model uses synthetically generated transaction data rather
than real NRB records. While distributions are grounded in
Nepal's context, results may differ with real data having
different statistical properties.

The model captures pairwise interactions only. In datasets with
strong multi-way correlations the model may underestimate total
risk.

The adversary model covers passive observation only. Active
attacks are outside scope.

Several Hyperledger Fabric-specific metadata fields are not
modelled --- transaction ID linkability, endorsing peer identity,
and certificate metadata. Including these would increase $m$ and
likely increase $\Rnorm$ for all schemes.

% ================================================================
\section{Future Work}
\label{sec:future_work}
% ================================================================

\textbf{Higher-order interactions.} Extending the model to
triplet and higher-order interaction terms would increase
measurement precision for datasets with strong multi-way
correlations, generalising to $\binom{m}{k}$ terms for order $k$.

\textbf{Real transaction data.} Applying the model to real data
from eSewa, Khalti, or NRB's own records would validate the
synthetic distributions and produce deployment-specific estimates
directly actionable for NRB.

\textbf{Formal connection to existing frameworks.} A formal proof
connecting $\Rnorm$ to the k-anonymity framework of
\citet{sweeney2002} and the differential privacy framework of
\citet{dwork2006} would position the metric within the
established theoretical privacy literature.

\textbf{Additional Fabric metadata attributes.} Extending to
include transaction ID linkability, endorsing peer identity, and
certificate metadata would raise $m$ from 7 to approximately 10,
producing a more complete risk assessment for real Fabric
deployments.

\textbf{Cross-border remittance extension.} Including remittance
corridor in a cross-border CBDC model would address Nepal's
significant international financial flows excluded from this
domestic-scope analysis.

% ================================================================
\section{Closing Remarks}
\label{sec:closing}
% ================================================================

Nepal is at a consequential moment in its financial history.
Millions of citizens who have never held a bank account may soon
transact digitally for the first time. That opportunity is real,
but so is the risk that the systems built to serve them could
expose their lives in ways they never consented to.

The choice of cryptographic scheme in a CBDC is not a technical
detail. It is a policy decision with direct consequences for
every person who uses the system. The difference between Plain
Hash Token and Ring Signature is a factor of 4.6 in
re-identification risk at national scale. The difference between
Ring Signature with full mitigation and no mitigation is a
factor of 34.

Privacy in a CBDC is not incompatible with financial inclusion,
regulatory compliance, or operational efficiency. Ring Signature
with a full mitigation stack delivers near-ZKP privacy at a
fraction of ZKP's computational cost. Nepal does not have to
choose between reaching its unbanked population and protecting
that population's privacy. It has to choose the right
cryptographic scheme.